\documentclass[12pt]{article}
\usepackage[T1]{fontenc}
\usepackage[utf8]{inputenc}
\usepackage[brazil]{babel}
\usepackage[margin=1in]{geometry}
\usepackage{ragged2e}
\usepackage[normalem]{ulem}
\setlength{\parindent}{0pt}
\setlength{\parskip}{0.8\baselineskip}
\begin{document}
\justifying
\noindent Verifica-se que a questão objeto da controvérsia jurídica suscitada no recurso manejado pela parte recorrente esteve afetada no representativo da controvérsia - \textbf{Tema n. }\textbf{164} da TNU - (PU no processo n. PEDILEF 0500774-49.2016.4.05.8305/PE),  afetado em 24/02/2017 foi julgado (trânsito em julgado em 02/10/2018), por meio do qual foi elucidada a seguinte questão:

\noindent \textit{"Saber quais são os reflexos das novas regras constantes na MP nº 739/2016 (§§ 8º e 9º do art. 60 da Lei 8.213/1991) na fixação da data de cessação do benefício auxílio-doença e da exigência, quando for o caso, do pedido de prorrogação, bem como se são aplicáveis aos benefícios concedidos e às demandas ajuizadas em momento anterior à sua vigência."}

\noindent Restou firmada a seguinte tese:

\noindent \textit{"Por não vislumbrar ilegalidade na fixação de data estimada para a cessação do auxílio-doença, ou mesmo na convocação do segurado para nova avaliação da persistência das condições que levaram à concessão do benefício na via judicial, a Turma Nacional de Uniformização, por unanimidade, firmou as seguintes teses: a) os benefícios de auxílio-doença concedidos judicial ou administrativamente, sem Data de Cessação de Benefício (DCB), ainda que anteriormente à edição da MP nº 739/2016, podem ser objeto de revisão administrativa, na forma e prazos previstos em lei e demais normas que regulamentam a matéria, por meio de prévia convocação dos segurados pelo INSS, para avaliar se persistem os motivos de concessão do benefício; b) os benefícios concedidos, reativados ou prorrogados posteriormente à publicação da MP nº 767/2017, convertida na Lei n.º 13.457/17, devem, nos termos da lei, ter a sua DCB fixada, sendo desnecessária, nesses casos, a realização de nova perícia para a cessação do benefício; c) em qualquer caso, o segurado poderá pedir a prorrogação do benefício, com garantia de pagamento até a realização da perícia médica." (Tema 164 TNU).}

\noindent Com efeito, depreende-se da leitura do voto condutor do acórdão que o mesmo se encontra em conformidade com o decidido no  tema pela TNU.

\noindent A proposito, consta do acordao hostilizado: (...) Rural (art. 59/63), Auxílio-Doença Previdenciário, Benefícios em Espécie, DIREITO PREVIDENCIÁRIO, Concessão, Pedidos Genéricos Relativos aos Benefícios em Espécie, DIREITO PREVIDENCIÁRIO, Parcelas de benefício não pagas, Pedidos Genéricos Relativos aos Benefícios em Espécie, DIREITO PREVIDENCIÁRIO

\noindent Nessas condições, o conteúdo do Acórdão recorrido é consentâneo com o entendimento da TNU, firmado em julgamento de recurso representativo de controvérsia, o que atrai a incidência da regra prevista pelo art. 14, III, \textit{b,} da Resolução CJF n. 586/2019:

\noindent \textit{\textbf{Art. 14.}}\textit{ Decorrido o prazo para contrarrazões, os autos serão conclusos ao magistrado responsável pelo exame preliminar de admissibilidade, que deverá, de forma sucessiva: (...) }\textit{\textbf{III }}\textit{- negar seguimento a pedido de uniformização de interpretação de lei federal interposto contra acórdão que esteja em conformidade com entendimento consolidado: }\textit{\textbf{b)}}\textit{ em }\uline{\textit{\textbf{recurso representativo de controvérsia pela Turma Nacional de Uniformização}}}\textit{ ou em pedido de uniformização de interpretação de lei dirigido ao Superior Tribunal de Justiça;(...).}

\noindent Posto isso, com arrimo no \uline{\textbf{art. 14, III, }}\uline{\textit{\textbf{b}}}\textit{,} da Resolução CJF n. 586/2019 c/c art. 1.030, I, do CPC, \textbf{NEGO SEGUIMENTO }\textbf{a}o\textbf{ PU}\textbf{.}

\noindent Intimem-se. Aguarde-se o decurso do prazo recursal. Após, certifique-se o trânsito em julgado e devolva-se ao juízo de origem.
\end{document}
